\documentclass[11pt,a4paper]{article}

\usepackage[T1]{fontenc}
\usepackage[utf8]{inputenc}
\usepackage{lmodern}
\usepackage[a4paper,margin=26mm]{geometry}
\usepackage{amsmath,amssymb,amsthm,mathtools}
\usepackage{booktabs,array}
\usepackage{microtype}
\usepackage{xcolor}
\usepackage{enumitem}
\usepackage{url}
\usepackage[colorlinks=true,linkcolor=blue!55!black,citecolor=blue!55!black,
            urlcolor=blue!55!black]{hyperref}

\newtheorem{theorem}{Theorem}
\newtheorem{lemma}[theorem]{Lemma}
\newtheorem{proposition}[theorem]{Proposition}
\newtheorem{remark}[theorem]{Remark}
\newcommand{\bits}{\{0,1\}}
\newcommand{\xor}{\mathbin{\oplus}}
\newcommand{\Sign}{\mathsf{Sign}}
\newcommand{\Verify}{\mathsf{Verify}}
\newcommand{\KeyGen}{\mathsf{KeyGen}}
\newcommand{\pk}{\mathsf{pk}}
\newcommand{\sk}{\mathsf{sk}}
\newcommand{\Hmsg}{H_{\mathsf{msg}}}
\newcommand{\Hseed}{H_{\mathsf{seed}}}
\newcommand{\Hiv}{H_{\mathsf{iv}}}
\newcommand{\TreeExpand}{\mathsf{TreeExpand}}
\newcommand{\LeafExpand}{\mathsf{LeafExpand}}
\newcommand{\VoleExpand}{\mathsf{VoleExpand}}
\newcommand{\BAVC}{\mathsf{BAVC}}
\newcommand{\VOLE}{\mathsf{VOLE}}
\newcommand{\pos}{\mathsf{pos}}
\newcommand{\Path}{\mathsf{Path}}
\newcommand{\Gala}{\mathsf{Gala}}
\newcommand{\Good}{\mathsf{Good}}
\newcommand{\E}{\mathbb E}
\newcommand{\Prb}{\mathop{\rm Pr}}

\title{\bfseries Cross-Variant Seed-Tree Reuse in Galas:\\
Secret-Key Recovery and Fresh-Message Forgery\\
from Same-Key, Same-Message S/F Signatures}
\author{Martin Feussner\\
\small Selmer Center, University of Bergen\\
\small \texttt{martin.feussner@uib.no}}
\date{1 October 2026}

\begin{document}
\maketitle

\begin{center}
\fbox{\parbox{0.92\linewidth}{\small\textbf{Provenance and review status.}
The attack, derivation, implementation, experiments, and initial manuscript were
produced by OpenAI Codex using Daybreak Blue at Ultra reasoning effort during an
AI-run audit of NGCC signature candidates, without a human first proposing this
cross-variant mechanism.  At the date above, no human cryptanalyst has independently
verified the argument, implementation, measurements, or novelty assessment.  This
report is a preliminary disclosure for human review and reproduction.}}
\end{center}
\vspace{0.7em}

\begin{abstract}
Galas supplies a short S variant and a fast F variant at each of four field sizes.
At a fixed field size the variants have the same untagged public and secret key
formats, the same algebraic signing relation, and deterministic signing.  We show a
practical specification-level cross-variant key/domain-separation defect when the
same valid key signs the same adversary-chosen message once under S and once under F.
The
seed derivation contains no variant identifier, so both signatures use the same
GGM root and IV.  The F opening hides one seed in its logical tree~0.  That F leaf
position is an internal node of the much larger S tree and, with high probability,
the S opening permits reconstruction of that node seed.  The two openings therefore
permit reconstruction of every F tree-0 seed.
The attacker recomputes the F witness mask $u_F$ and uses the serialized relation
\[
             w=d_F\xor u_F[0..\ell),\qquad k=w[0..\lambda)
\]
to recover the exact secret key.  No hash inversion, algebraic key search, or
Fiat--Shamir forgery is needed.

At each field size, the attack recovered the exact key from one same-key,
same-message pair of signatures emitted by the untouched submitted S/F interfaces.
After converting the recovered normative $k$ only to the wrapper's documented
$x\mathbin\|k$ cache encoding, the untouched F signer produced a fresh-message
signature accepted by the untouched F verifier.  Offline recovery took 0.05,
0.25, 0.60, and 1.19 seconds and at most 24 MiB on an AMD EPYC host.  For matched
valid inputs, the original and repaired signers produced byte-identical complete
serialized transcripts in all eight S/F profiles.  A separately written Galas-160
implementation recovered and forged on five further keys; a 512-message sweep
under the real signing distribution gave 509 complementary openings.  An ideal
unconditioned index model predicts success
probabilities between 0.9952 and 0.9975, but these are not exact serialized-signature
probabilities because grinding, opening-cap conditioning, and the shared root can
correlate the challenges.  Failed pairs are publicly detectable and can be retried
with another common message.

The result is conditional on cross-variant key reuse and access to both signing
interfaces on a common message.  It does not break the ordinary isolated,
single-variant EUF-CMA game, and it does not show that deployments which enforce
separate typed keys are vulnerable.  The specification does not instruct users to
reuse a key across S and F.  It also neither prohibits that reuse nor
cryptographically separates the S/F key namespaces: $\KeyGen(1^\lambda)$ has no
variant input and $\sk=k$ has no variant tag.  Binding the full instance identifier
into seed derivation and the key namespace removes the attack.
\end{abstract}

\section{Introduction}

Galas is a VOLE-in-the-Head signature built around the Gala keyed one-way
function~\cite{galas-spec,galas-design}.  At each submitted field size
$\lambda\in\{160,256,384,512\}$, it offers an S variant intended to reduce
signature size and an F variant intended to improve throughput.  The variants
change the one-tree BAVC/VOLE schedule while retaining the same Gala relation.

This flexibility interacts badly with deterministic signing.  The normative key
and seed derivation are parameterized by $\lambda$, the public key, the secret key,
and the message, but not by the S/F label.  Consequently, same-key, same-message
signatures from the two variants use the same root seed in two differently sized
GGM trees.  A node that is a hidden leaf in the F tree is an ordinary internal node
in the S tree.  The S opening normally permits reconstruction of that node seed.
Reinterpreting its seed
under the public F leaf-expansion rule completes the F sender's first VOLE vector,
which is precisely the one-time pad applied to the signing witness.

Our contributions are as follows.
\begin{enumerate}[leftmargin=*,itemsep=2pt]
  \item We prove the exact shared-node identity and the complementary-opening
        condition directly from the normative algorithms.
  \item We give a two-query-per-attempt recovery algorithm whose success or
        failure can be detected using only public transcript data.  On success it
        recovers the original $\lambda$-bit key, rather than an equivalent
        representation.
  \item We demonstrate exact key recovery from original-interface transcripts and
        fresh-message forgery accepted by the untouched submitted verifier at all
        four field sizes.  A separately written hostile reproducer confirms the
        Galas-160 mechanism and matched negative controls.
  \item We quantify the seed-tree geometry, while separating a rigorous
        unconditioned index calculation from the conditioned real-signer
        distribution.  A 512-message experiment measures the latter directly.
  \item We delimit the security claim: this is a cross-variant attack under a
        same-key, same-message prerequisite, not a one-oracle forgery against an
        isolated variant.
\end{enumerate}

The finding is a specification-level cross-variant key/domain-separation defect
under same-key reuse.  The specification does not require that reuse.  We tested
both a repaired specification oracle and the untouched submitted interfaces.  For
the ordinary valid inputs used here, all eight original S/F signers produced the
same complete serialized transcripts as the oracle, and all fresh forgeries were
accepted by the untouched submitted verifiers.  The attack does not use the
public-zero key-generation
fallback or message-length truncation already reported on the NGCC site.

\section{Relevant construction}

\subsection{Keys, witness, and deterministic signing}

For a fixed even field size $\lambda$, Galas works over
$\mathbb F_{2^\lambda}$.  Algorithm~19 of the specification samples $k,x$, tests
$\Good(k,x)$, and returns
\begin{equation}
  \pk=x\mathbin\|y,\qquad \sk=k,\qquad y=\Gala_k(x).       \label{eq:keygen}
\end{equation}
The algorithm is written as $\KeyGen(1^\lambda)$: it has no S/F input.  At a
fixed $\lambda$, both variants use the same $2\lambda$-bit public key, the same
$\lambda$-bit secret key, and the same witness
\[
  w=(k,\sigma_0,\sigma_1,\sigma_2)\in\bits^\ell,
  \qquad \ell=5\lambda/2.
\]
The $\sigma_i$ are deterministic inverse-norm witnesses computed from $(\pk,k)$.

Algorithm~21 is deterministic.  Given the same $(\pk,k,M)$, it first computes
\begin{equation}
 \mu=\Hmsg(\pk,M),\qquad
 (K_{\rm root},iv_{\rm pre})=\Hseed(k,\mu),\qquad
 iv=\Hiv(iv_{\rm pre}).                                 \label{eq:state}
\end{equation}
No instance name is inserted in these calls.  The BAVC/VOLE layer produces a
mask $u$, and the signature exposes
\begin{equation}
                   d=w\xor u[0..\ell).                  \label{eq:maskedwitness}
\end{equation}
The complete signature also contains the BAVC opening, final challenge $\Delta$,
$iv_{\rm pre}$, and the grinding counter.  Verification reconstructs the opened
VOLE transcript and checks the final Fiat--Shamir equation, without learning
$w$ or $u$.

\subsection{One-tree BAVC geometry}

For a concrete variant $v\in\{S,F\}$, let $\tau_v$ be the number of logical
VOLE vectors.  Logical vector $i$ has $N_{v,i}=2^{k_{v,i}}$ leaves.  All logical
leaves are interleaved in one binary heap with
\[
                  L_v=\sum_{i=0}^{\tau_v-1}N_{v,i}
\]
physical leaves and $2L_v-1$ nodes.  The root is node~0 and node $\alpha$ has
children $2\alpha+1$ and $2\alpha+2$.  We write $\pos_v(i,j)$ for the normative
physical position of logical leaf $(i,j)$.

Starting at $K_0=K_{\rm root}$, every internal node expands as
\begin{equation}
 K_{2\alpha+1}\mathbin\|K_{2\alpha+2}
   =\TreeExpand(K_\alpha,iv,\alpha),                    \label{eq:treeexpand}
\end{equation}
where the PRG frame is
\[
 K_\alpha\mathbin\|iv\mathbin\|\mathsf{u32le}(\alpha)
          \mathbin\|\mathsf{u32le}(\mathit{counter}).
\]
It contains neither the variant label nor $L_v$.  A physical leaf seed is then
expanded using the variant's public leaf tweak:
\begin{equation}
 sd_{v,i,j}\mathbin\|com_{v,i,j}
 =\LeafExpand(K_{\pos_v(i,j)},iv,L_v-1+i).              \label{eq:leafexpand}
\end{equation}

The final challenge selects one hidden index $J_v[i]$ in each logical vector.
Define the union of hidden paths
\begin{equation}
 \mathcal H_v=\bigcup_{i=0}^{\tau_v-1}
   \Path\bigl(0,\pos_v(i,J_v[i])\bigr).                 \label{eq:hiddenpaths}
\end{equation}
The BAVC opening provides sibling seeds from which every node outside
$\mathcal H_v$ is reconstructed.  In particular, it permits reconstruction of all
logical-tree-0 leaf seeds except $K_{\pos_v(0,J_v[0])}$.

Finally, the first logical vector is converted to the sender mask by
\begin{align}
 R_{0,j}&=\VoleExpand(sd_{v,0,j},iv,0),\nonumber\\
 R_{r+1,j}&=R_{r,2j}\xor R_{r,2j+1},\qquad
 u_v=R_{k_{v,0},0}.                                    \label{eq:volefold}
\end{align}
The complete VOLE commitment sets $u=u_v$ from logical vector~0.

\begin{table}[t]
\centering
\caption{Normative S/F schedules and serialized signature sizes.}
\label{tab:params}
\begin{tabular}{lrrrrrr}
\toprule
Instance & $\lambda$ & $\tau$ & $T_{\rm open}$ & $k_{\rm tree}$ & $L$ & bytes\\
\midrule
Galas-160S & 160 & 14 & 132 & 11 & 27,648 & 4,812\\
Galas-160F & 160 & 20 & 144 &  8 &  4,096 & 5,964\\
Galas-256S & 256 & 21 & 218 & 12 & 79,872 & 12,114\\
Galas-256F & 256 & 35 & 232 &  7 &  4,416 & 15,950\\
Galas-384S & 384 & 32 & 336 & 12 &118,784 & 27,784\\
Galas-384F & 384 & 48 & 332 &  8 & 11,520 & 33,384\\
Galas-512S & 512 & 42 & 456 & 13 &172,032 & 49,516\\
Galas-512F & 512 & 64 & 445 &  8 & 15,616 & 59,416\\
\bottomrule
\end{tabular}
\end{table}

\section{The cross-variant identity}

\begin{lemma}[Shared deterministic state]\label{lem:state}
Fix $\lambda$, a valid normative key pair $(\pk,k)$, and a message $M$.  An S
signature and an F signature generated from these values have identical
$\mu,K_{\rm root},iv_{\rm pre}$, and $iv$.
\end{lemma}
\begin{proof}
Every quantity in~\eqref{eq:state} is a deterministic function of the displayed
arguments.  The specification inserts no S/F label or full instance identifier
in $\Hmsg$, $\Hseed$, or $\Hiv$.  The inputs are therefore byte-for-byte equal.
\end{proof}

\begin{lemma}[Shared physical nodes]\label{lem:nodes}
Under the conditions of Lemma~\ref{lem:state}, a physical GGM node index
$\alpha$ that is generated in both heaps has the same seed in the S and F
executions.
\end{lemma}
\begin{proof}
Both heaps start with the same $K_0$ and $iv$.  Equation~\eqref{eq:treeexpand}
is a deterministic function of $(K_\alpha,iv,\alpha)$ and contains no schedule
identifier.  Induction down the common heap indices proves the claim.
\end{proof}

The leaf tweak in~\eqref{eq:leafexpand} does contain $L_v$, but this does not
separate the underlying node seeds.  Once an attacker reconstructs $K_\alpha$
from the S opening, the attacker can apply the public F tweak to it.

\begin{proposition}[Every F leaf is an internal S node]\label{prop:heapbound}
For each of the four submitted S/F pairs, every physical F leaf position---and
therefore every possible F logical-tree-0 target $\alpha^*$---is the index of an
internal node in the corresponding S heap.
\end{proposition}
\begin{proof}
A binary heap with $L$ leaves numbers its internal nodes from $0$ through $L-2$
and its leaves from $L-1$ through $2L-2$.  Thus every F leaf index is at most
$2L_F-2$, whereas an index is internal in the S heap whenever it is strictly
less than $L_S-1$.  The required strict inequality holds at every field size:
\begin{center}
\begin{tabular}{lrrc}
\toprule
Pair & $2L_F-2$ & $L_S-1$ & $2L_F-2<L_S-1$\\
\midrule
160S/F &  8,190 &  27,647 & yes\\
256S/F &  8,830 &  79,871 & yes\\
384S/F & 23,038 & 118,783 & yes\\
512S/F & 31,230 & 172,031 & yes\\
\bottomrule
\end{tabular}
\end{center}
The bound covers all F leaves, so no additional property of the interleaving map
or of the selected tree-0 index is needed.
\end{proof}

\begin{theorem}[Conditional exact key recovery]\label{thm:recovery}
Let
\[
             \alpha^*=\pos_F(0,J_F[0]).
\]
Suppose $\alpha^*\notin\mathcal H_S$.  From two valid, honestly generated
same-key, same-message S/F signatures and the public parameters, an attacker
recovers the exact normative secret key $k$ in linear PRG work in the two heap
sizes.
\end{theorem}
\begin{proof}
The F opening reconstructs every logical-tree-0 physical seed except
$K_{\alpha^*}$.  Proposition~\ref{prop:heapbound} shows that $\alpha^*$ is an
internal node of the S heap for every submitted pair.  Because
$\alpha^*\notin\mathcal H_S$, the S opening permits reconstruction of its seed.
By
Lemma~\ref{lem:nodes}, this is the same seed required by F.

The attacker now has all $K_{\pos_F(0,j)}$.  Applying the public F leaf expansion
in~\eqref{eq:leafexpand}, followed by the public fold in~\eqref{eq:volefold},
gives the exact $u_F$.  The F signature contains $d_F$, so
\begin{equation}
       w=d_F\xor u_F[0..\ell),\qquad k=w[0..\lambda).   \label{eq:recover}
\end{equation}
The first $\lambda$ witness bits are exactly the normative secret key.  The
algorithm uses arrays linear in $L_S+L_F$ and a constant number of public PRG
expansions per reconstructed node or F tree-0 leaf.
\end{proof}

The theorem explains the whole attack chain:
\[
\boxed{\begin{array}{c}
\text{missing S/F separation}\ \Longrightarrow\ \text{shared GGM nodes}\\
\Longrightarrow\ \text{complementary BAVC opening}
\Longrightarrow\ u_F\Longrightarrow w\Longrightarrow k
\Longrightarrow\text{fresh-message signing.}
\end{array}}
\]
No property of the Gala one-way function is weakened.  Instead, combining the
two openings permits reconstruction of a mask that the proof system relies on
keeping hidden.

\section{Success probability and retry strategy}

For a fixed F target node $\alpha$, define
\[
 c_i(\alpha)=\#\{j:\alpha\text{ is an ancestor of }\pos_S(i,j)\}.
\]
If the S indices $J_S[i]$ are independent and uniform before grinding and
opening-cap rejection, the probability that no S hidden path contains $\alpha$
is exactly
\begin{equation}
 P_\alpha=\prod_{i=0}^{\tau_S-1}
                 \left(1-\frac{c_i(\alpha)}{N_{S,i}}\right). \label{eq:palpha}
\end{equation}
If the F tree-0 index is also independent and uniform, averaging gives
\begin{equation}
 P_{S,F}=\frac1{N_{F,0}}
   \sum_{j=0}^{N_{F,0}-1}P_{\pos_F(0,j)}.              \label{eq:pairprob}
\end{equation}
An exact enumeration of the normative schedules gives Table~\ref{tab:prob}.

\begin{table}[ht]
\centering
\caption{Unconditioned independent-index model.  These are not claimed as exact
probabilities for serialized signatures.}
\label{tab:prob}
\begin{tabular}{lrrr}
\toprule
Pair & $N_{F,0}$ & $P_{S,F}$ & expected pairs $1/P_{S,F}$\\
\midrule
160S/F & 256 & 0.996787101321 & 1.00322\\
256S/F & 128 & 0.995242959290 & 1.00478\\
384S/F & 256 & 0.997515982924 & 1.00249\\
512S/F & 256 & 0.997348606864 & 1.00266\\
\bottomrule
\end{tabular}
\end{table}

The qualification matters.  A real signer accepts only final challenges that
satisfy the grinding predicate and whose canonical BAVC opening uses at most
$T_{\rm open}$ sibling seeds.  The S and F challenges also arise from transcripts
with a common root.  Equation~\eqref{eq:pairprob} therefore describes the
unconditioned schedule geometry, not a proof of the real success probability.

The real condition is nonetheless publicly testable.  After reconstructing both
openings, the attacker checks whether the S opening permits reconstruction of
$K_{\alpha^*}$.  If it does not, the pair is discarded before any secret-dependent
test.  The attacker asks
both interfaces to sign a new common chosen message and repeats.  Under the
random-oracle model, the distinct input changes $\mu$ and the deterministic root
except with negligible collision probability.

To measure the conditioned implementation rather than infer it from
Equation~\eqref{eq:pairprob}, we generated 512 same-key S/F signature pairs on
independently derived common messages.  Recovery succeeded on 509 and gave a
clean ``unrecovered target leaf'' exit on three, for an observed rate
$509/512=0.994140625$.  The sweep used 48 workers and completed in 34.10 seconds
of wall time on the audit host.  The experiment used one key, so it measures
message-to-message behavior for that key and is not a 512-key estimate.

\section{Recovery algorithm and public validation}

The attack on one pair is the following.
\begin{enumerate}[leftmargin=*,itemsep=2pt]
  \item Query $\Sign_S(\pk,k,M)$ and $\Sign_F(\pk,k,M)$ for the same chosen
        $M$, and verify both public signatures.
  \item Parse $iv_{\rm pre}$ and $\Delta$ from each signature.  Recompute $iv$,
        decode the opening indices, and reconstruct every node seed permitted by
        each canonical BAVC opening.
  \item Locate $\alpha^*=\pos_F(0,J_F[0])$.  Use the F view for the other
        $N_{F,0}-1$ physical seeds and reconstruct $K_{\alpha^*}$ from the S
        opening.  If the S opening does not permit that reconstruction, request
        another common-message pair.
  \item Apply $\LeafExpand$ with the F tweak $L_F-1$ to every recovered
        tree-0 seed.  Apply $\VoleExpand$ and the public XOR fold to obtain $u_F$.
  \item Compute $w$ and $k$ with~\eqref{eq:recover}.
  \item Validate the candidate without the retained secret by checking
        $\Good(k,x)=1$ and $\Gala_k(x)=y$ for the public $\pk=x\mathbin\|y$.
        A valid candidate is then an ordinary signing key for either variant.
\end{enumerate}

At $\lambda=160$, the two queries return $4{,}812+5{,}964=10{,}776$ bytes.
The two main node arrays contain
\[
 (2\cdot27{,}648-1)\cdot20+(2\cdot4{,}096-1)\cdot20
   =1{,}269{,}720\ \text{bytes},
\]
with small additional bitsets and work buffers.  The algorithm is deterministic
after it receives a successful pair.

Fresh-message forgery requires no special construction: run the ordinary
normative signing algorithm with the recovered $k$ on a message that was not
queried.  Verification acceptance is therefore both the end-to-end impact and an
independent public check that the recovered string is a functioning secret key.

\section{Specification-conformant and pristine-source experiments}

\subsection{Interface repairs and untouched-source equivalence}

The official specification PDF, rather than the submitted wrapper, defines the
design.  The companion \texttt{reproducer/README.md} records checksums for the
official specification and submission archive.  The reproducer includes a
specification-conformant
wrapper derived from the submitted sources.  Three interface differences were
repaired before testing:
\begin{enumerate}[leftmargin=*,itemsep=2pt]
  \item key generation requires an explicit independent seed instead of using
        the wrapper's public all-zero fallback;
  \item the normative secret key is $k$, while the source wrapper caches
        $x\mathbin\|k$ because its external signing API omits $\pk$; and
  \item the full host-size message length is passed to $\Hmsg$, instead of the
        source wrapper's 32-bit length.
\end{enumerate}
The repaired entry points implement Algorithms 19, 21, and 22 as
$\KeyGen$, $\Sign(\pk,k,M)$, and $\Verify(\pk,M,\sigma)$.  We then extracted the
official archive afresh and built all eight submitted profiles without
editing their source files.  For the same explicit seed and valid inputs, the
normative key $k$ equalled the suffix of the source wrapper's cached
$x\mathbin\|k$, while the prefix equalled the public $x$.  For every S/F profile,
the repaired and original signers emitted byte-identical complete serialized
signatures.  Each untouched verifier accepted both transcripts, and each repaired
verifier accepted the untouched-source transcript.  Table~\ref{tab:transcripts}
summarizes these checks.

\begin{table}[ht]
\centering
\caption{Whole-transcript comparison on matched valid inputs.  ``Cross-accept''
means that the untouched and repaired verifiers both accepted both serialized
transcripts.}
\label{tab:transcripts}
\begin{tabular}{lrrcc}
\toprule
Profile & bytes & SHA-256 prefix & exact equality & cross-accept\\
\midrule
Galas-160S &  4,812 & \texttt{9859a1242bb7d7b5} & yes & yes\\
Galas-160F &  5,964 & \texttt{0cacfff1200e847e} & yes & yes\\
Galas-256S & 12,114 & \texttt{7d73c06255b34650} & yes & yes\\
Galas-256F & 15,950 & \texttt{d30ee52a8b2276df} & yes & yes\\
Galas-384S & 27,784 & \texttt{166e2bc1bc1cdc86} & yes & yes\\
Galas-384F & 33,384 & \texttt{c27ae298c8b8eecd} & yes & yes\\
Galas-512S & 49,516 & \texttt{fda7ff0f8f5cf607} & yes & yes\\
Galas-512F & 59,416 & \texttt{2d5115b086ec3c26} & yes & yes\\
\bottomrule
\end{tabular}
\end{table}

The pristine experiment hashed all 552 files in the extracted reference-source
tree before and after building and testing; every hash remained unchanged.  Its
file-I/O harness calls the submitted \texttt{sig\_keygen}, \texttt{sig\_sign}, and
\texttt{sig\_verify} functions directly.  After recovery, its only conversion is
to place the normative $k$ into the wrapper's documented $x\mathbin\|k$ cache,
using public $x$, before calling the untouched signer.  Thus neither recovery nor
verification depends materially on the repaired wrapper, and the defect is not
created by any repaired deviation.

For portable reproduction, the repository reproducer also replaces the
submitted SM3 rotate macro's undefined shift by 32 when its rotation count is
zero with an equivalent defined C99 operation.  This cleanup leaves the fixed
seed transcripts unchanged under GCC.  An AddressSanitizer/UndefinedBehaviorSanitizer
run of the complete 160-bit recovery, forgery, and negative-control path then
completed without a diagnostic.  The pristine build retained the submitted macro
unchanged.

\subsection{All-level results}

Table~\ref{tab:results} reports the pristine-source run at every submitted field
size.  Each row used one S signature and one F signature emitted by the original
submitted interfaces under the same normative key on a common chosen message.
The recovery program received
only those two serialized signatures.  It did not receive the public key, cached
secret, or normative key.  After recovery, the experiment compared the result to
the independently retained key, signed a fresh message through the untouched F
signer, and called the untouched F verifier.  ``Wrong message'' applies that
fresh signature to a different message.

\begin{table}[ht]
\centering
\caption{End-to-end recovery through the original submitted interfaces.  Times
and RSS are for offline recovery, excluding signing-oracle time.}
\label{tab:results}
\begin{tabular}{lrrrrcc}
\toprule
Pair & query bytes & key bytes & time (s) & RSS (MiB) & fresh F & wrong message\\
\midrule
160S/F & 10,776  & 20 & 0.05 &  2 & accept & reject\\
256S/F & 28,064  & 32 & 0.25 &  6 & accept & reject\\
384S/F & 61,168  & 48 & 0.60 & 13 & accept & reject\\
512S/F &108,932  & 64 & 1.19 & 24 & accept & reject\\
\bottomrule
\end{tabular}
\end{table}

Every recovered key matched both the normative $k$ suffix of the untouched
wrapper secret and the repaired oracle's Algorithm-19 key byte for byte.  The
160-bit retained and recovered key had SHA-256
\begin{center}
\small\texttt{01f62f9bd378ef93973c4c172ba0b804377b165244c4eb6d66319ac24e50bb76}.
\end{center}
The 160-bit fresh-message F forgery generated through the untouched signer had
SHA-256
\begin{center}
\small\texttt{3de6b380fa1b66f3b8e9c6f7afb68e735affdacc6bd8fc1fc087aed75fcbf952}
\end{center}
and the untouched verifier accepted it.  The untouched verifier took 0.77,
3.12, 11.06, and 33.25 seconds respectively.

\subsection{Independent AI reproduction and controls}

A separately instructed Codex agent wrote a new 160-bit recovery implementation
against a clean extraction of the official archive.  It did not reuse the initial
recovery source.  Across five independently seeded keys it obtained five exact key
matches and five accepted fresh-message forgeries.  A timed recovery took 0.06
seconds and 2,048 KiB RSS.  A symbol and file-access audit found no linked
key-generation, signing, OWF, retained-secret, or instrumentation input in the
recovery executable.  AddressSanitizer reported no error on the successful path.

The following adversarial controls had their expected outcomes:
\begin{itemize}[leftmargin=*,itemsep=2pt]
  \item the F opening alone stopped at its one hidden tree-0 seed;
  \item pairing different keys or pairing different messages failed the shared
        state check and emitted no key;
  \item swapping the S and F signature arguments failed fixed-length parsing;
  \item the original chosen-message signature and the fresh forgery were each
        rejected when verified on a different message;
  \item each fresh forgery was rejected under an independently generated public
        key, and a signature made with an independently generated wrong secret
        was rejected under the target public key;
  \item one-bit signature mutations were rejected by Algorithm~22;
  \item changing the serialized $d_F$ changed the recovered string, which then
        failed to produce an accepted fresh signature; and
  \item in that independent 160-bit experiment, signatures produced by the
        original source and repaired oracle were byte-identical for both variants.
\end{itemize}

Separately, the pristine-source validation established whole-transcript equality
for all eight S/F profiles and the all-level untouched-verifier results in
Tables~\ref{tab:transcripts} and~\ref{tab:results}.  This separates the original
separately written 160-bit implementation from the later all-profile validation.

This second implementation was also produced within the AI audit and is not
independent human verification.  Its results confirm the practical
specification-level cross-variant key/domain-separation defect under same-key
reuse; they do not establish an isolated single-variant EUF-CMA break.

\subsection{Reproduction}

The companion repository contains a deterministic reproducer.  From the
\texttt{reproducer} directory, run
\begin{center}
\texttt{./run.sh}.
\end{center}
The command generates one Galas-160 key, obtains one S signature and one F
signature on the same chosen message, recovers the exact key from the two
serialized signatures, and produces a fresh-message F signature.  It also
requires rejection of wrong-message, wrong-key, mutated-signature, and
nonmatching-transcript controls.  A successful run ends with
\texttt{minimal\_reproducer=pass}.

The optional command \texttt{./run\_all\_levels.sh} repeats exact recovery and
fresh-message forgery at 256, 384, and 512 bits and ends with
\texttt{all\_levels\_reproducer=pass}.  The file
\texttt{reproducer/README.md} records the required tools, expected running time,
source-archive checksum, and the separate untouched-source validation procedure.
The recovery program receives only the two serialized signatures; comparison
with the retained original key is a post-recovery experiment control.

\section{Security impact and limits}

The demonstrated impact is exact recovery of $k$, followed by an ordinary
fresh-message signature accepted by the normative verifier.  It is stronger than
transcript malleability and is not merely a reduction in a security estimate.
It also recovers the original key, not only a proof-system witness or equivalent
key.  The prerequisites and nonclaims are equally concrete.

\paragraph{Cross-variant oracle access.}
The adversary needs a geometrically successful pair consisting of one S and one F
signature under the same public/secret pair $(\pk,k)$ and on the same message.
Each attempt uses two signing queries, and the pair's geometric success is
publicly testable.  The ordinary EUF-CMA game for one fixed algorithm instance
offers only one signing oracle.  This attack therefore does not break Galas-160S
alone, Galas-160F alone, or any other isolated variant in that strict game.

\paragraph{Classification and same-key prerequisite.}
This is a specification-level cross-variant key/domain-separation defect under
same-key reuse.  The specification does not explicitly instruct users to reuse
keys across S and F.  It defines $\KeyGen(1^\lambda)$ without an S/F input, gives
both variants the identical encodings in~\eqref{eq:keygen}, states that the
algebraic relation is the same, and defines no key-usage prohibition.  The
repaired and submitted core signers accept the same serialized key material in
both variants.  Cross-variant reuse is therefore neither prohibited nor
cryptographically separated by the specified algorithms, key namespace, or
serialization.  An API or deployment that generates and enforces a separate
typed key for each variant avoids the prerequisite.

\paragraph{Same-message requirement.}
Different messages change $\mu$ and hence the root.  Repeated signing of one
message within a single variant returns the same deterministic transcript and
hides the same paths, so it does not provide complementary information.  The
attack neither claims nor needs a same-variant nonce-reuse weakness.

\paragraph{Success estimates.}
Table~\ref{tab:prob} is a rigorous calculation only for independent uniform
decoded indices before the signer's grinding and $T_{\rm open}$ conditioning.
The $509/512$ experiment supports practical success for one 160-bit key but does
not prove the all-key probability or exact independence between retries.  A failed
pair does not yield a key in this procedure and is handled by a publicly selected
retry.

\paragraph{No primitive cryptanalysis.}
The recovery relies on exact equality of PRG inputs, public BAVC reconstruction,
and the witness-mask equation.  It assumes neither a collision nor a weakness in
the NGCC hash/XOF, the Gala relation, QuickSilver, or the Fiat--Shamir transform.
The already public zero-seed and message-length source findings are unrelated.

\section{Repair}

The root cause is a missing full-instance domain separator at the point where
long-lived secret material derives the deterministic commitment state.  A direct
repair is
\begin{equation}
 (K_{\rm root},iv_{\rm pre})=
 \Hseed(\mathsf{instanceID}\mathbin\|k\mathbin\|\mu),   \label{eq:repair}
\end{equation}
where \textsf{instanceID} canonically includes the field size and S/F label.
Binding the same identifier into $\Hmsg$ and every tree, leaf, and VOLE expansion
frame provides defense in depth and makes cross-protocol equality visibly
impossible at each layer.

Keys and APIs should also be typed by the complete instance.  A serialized key
created for Galas-160S should be rejected by a Galas-160F signing entry point even
though the algebraic relation has the same shape.  Known-answer tests should
assert that equal raw $(\pk,k,M)$ inputs under S and F produce different roots and
IVs.  Either independent roots or enforced disjoint key namespaces blocks the
attack; deploying both is safer and easier to audit.

\section{Prior art and novelty boundary}

The public NGCC Galas page listed two critical implementation reports at the
1~October~2026 checkpoint: a public-zero key-generation fallback and 32-bit
message-length truncation~\cite{ngcc-galas}.  Neither report concerns shared
S/F roots, complementary BAVC openings, or recovery of $u$ and $k$.

Chosen-protocol attacks established the broader warning that a key may be secure
inside either of two protocols yet fail when both protocols expose the same key
material~\cite{kelsey1998}.  Modified-parameter attacks later made the closely
related point that security under one parameter set need not survive use of the
same key under another~\cite{howgrave2004}.  These works provide the general
cross-protocol and cross-parameter perspective.  They do not identify the Galas
shared-node identity, the complementary S/F openings, or recovery of the VOLE
witness mask; those are the Galas-specific claims here.

FAEST, the closest transcript template cited by Galas, derives its root and IV
from secret and message material together with a fresh $\lambda$-bit $\rho$ by
default; deterministic operation sets $\rho=0$~\cite{faestv2}.  This comparison
illustrates why fresh or domain-separated state prevents accidental equality, but
it is not prior art for the attack proved here.

Targeted searches covered the NGCC reports and security index, the local public
NGCC harness snapshot, the PKC Forum, the IACR ePrint Archive, the Galas design
paper title and authors, and combinations of ``Galas'', ``cross-variant'',
``same key'', ``seed tree'', and ``parameter attack''.  We found no public report
of complementary Galas S/F openings recovering the VOLE witness mask.  This is a
best-effort negative search as of 1~October~2026, not proof that unpublished or
unindexed prior art does not exist.  The live sources were last checked on
1~October~2026.

\section{Conclusion}

Two Galas signatures can expose an entire signing key when the same valid key
signs the same chosen message once under S and once under F.  Missing variant
separation makes their deterministic GGM nodes identical; the larger S opening
normally permits reconstruction of the one F tree-0 node seed hidden by the F
opening.  Public expansion then recovers $u_F$, the masked witness, and the exact
$k$.  Full-parameter
experiments recovered 20-, 32-, 48-, and 64-byte keys in at most 1.19 seconds and
used them to produce fresh-message signatures accepted by the untouched submitted
verifiers.

The precise scope should remain attached to the result.  It is a practical
cross-variant family attack under same-key, same-message access, while the
isolated single-variant EUF-CMA games remain unbroken by this method.  Full-instance
domain separation and typed keys remove the shared state that enables recovery.

\section*{Acknowledgements}

The computations were performed on the Norwegian Research and Education Cloud
(NREC), using resources provided by the University of Bergen and the University
of Oslo.

\begin{thebibliography}{9}
\small

\bibitem{galas-spec}
Guoxiao Liu, Zhuoqing Peng, Chongxu Ren, Guowei Liu, Juntong Lin, Yimeng Sun,
Keting Jia, and Hongbo Yu,
\emph{Galas Signature Scheme: Submission to ICCS NGCC}, 26 June 2026.

\bibitem{galas-design}
Guoxiao Liu, Yimeng Sun, Zhuoqing Peng, Chongxu Ren, Guowei Liu, Juntong Lin,
Keting Jia, and Meiqin Wang,
\emph{GALAS: Compact VOLE-in-the-Head Signatures from Subfield Compression},
ACM CCS 2026, accepted Cycle~1, to appear.

\bibitem{ngcc-galas}
NGCC Security Reports,
\emph{Galas (sign-12)},
\url{https://ngcc.dev/reports/sign-12.html}, accessed 1 October 2026.

\bibitem{kelsey1998}
John Kelsey, Bruce Schneier, and David Wagner,
\emph{Protocol Interactions and the Chosen Protocol Attack},
Security Protocols, 5th International Workshop, Lecture Notes in Computer
Science 1361, pp.~91--104, Springer, 1998,
\url{https://doi.org/10.1007/BFb0028162}.

\bibitem{howgrave2004}
Nick Howgrave-Graham, Joseph H. Silverman, Ari Singer, and William Whyte,
\emph{Modified Parameter Attacks: Practical Attacks Against CCA2 Secure
Cryptosystems, and Countermeasures},
IACR Cryptology ePrint Archive, Report 2004/344, 2004,
\url{https://eprint.iacr.org/2004/344}.

\bibitem{faestv2}
The FAEST Team,
\emph{FAEST Specification, Version 2.0},
\url{https://faest.info/faest-spec-v2.0.pdf}.

\end{thebibliography}

\end{document}
